\scanpage{61}
% archive sheet number: 60

\deleted{\deleted{et en choisir} \inserted{d'où une arête privilégiée sur $D$} (celle qui rencontre $D_0$), et je choisis une extrémité $s$, de façon à avoir un repère. Je spécialise $D$ en $D'$ vers le \deleted{pt} \inserted{sommet} $s'$ le plus proche de $s$, je fais un demi-tour par $s$, et je reviens à $D$, --- un nb \deleted{pair} d'opérations, qui change l'orientation de $D$, \uline{et transforme $s$ en $s$} --- donc c'est la symétrie p.r. à $s$. Si on \deleted{fait}}

4.1.

Pour tout \deleted{lacet} chemin \inserted{$c$} combinatoire sur $\Gamma$ --- i.e. chemin formé avec des faces $\unsure{\overline F_D|F_D}$ de $(\mathcal X,K)$, \deleted{\unsure{quel}}\unsure{reliant} la face $F=F_D$ \unsure{avec|une} $F'=\overline F_{D'}$, on trouve un iso ``de transport parallèle''

\[ \operatorname{Pol}(c):\quad \operatorname{Pol}(F)\unsure{\text{---}|\longrightarrow}\operatorname{Pol}(F'), \]

qui du pt de vue de la géométrie intrinsèque de $(\mathcal X,K,S_K)$ s'interprète par composition successive de ``réflexions'' le long des arêtes de réflexion qui interviennent au cours de $c$.

L'effet de $\operatorname{Pol}(c)$ sur les orientations

\[ \underline{\omega}_x(c):\quad \underline{\omega}(F)\simeq\underline{\omega}(D)\longrightarrow\underline{\omega}(F')\simeq\underline{\omega}(D') \]

est aussi l'effet sur l'orientation des p.r. par transport parallèle --- donc \deleted{est} indépendant de la parité du nb de pas dans le chemin $c$ ! Si on a un lacet \deleted{\unsure{\ell}} $l=c$ i.e. $F=F'$ --- donc

\[ \Pi_1(\mathcal X,F_D)\longrightarrow\underset{\simeq\mathbb D_{2n}}{\operatorname{Aut}(\operatorname{Pol}(D))}\xrightarrow{\substack{\text{effet sur}\\\text{orientation}}}\pm1 \]

on désignera par $\chi_x$ le caractère composé

\[ \chi_x:\Pi_1(\mathcal X,F_D)\longrightarrow\pm1 \]

donc pour un chemin $l$, on aura $\chi_{\unsure{,|x|\ast}}(l)=+1$ ou $-1$ suivant que par transport parallèle l'orientation de $D$ aura ou non changé. \deleted{\unsure{On}}

On a deux autres caractères $\chi_s$ et $\chi_a$ (\deleted{\unsure{orient}}ations ``sommets'' et ``arêtes'') en associant à chaque face $F_D$ \deleted{\unsure{les}} deux ensembles à deux éléments

\[ \begin{aligned}
\varepsilon_s(\deleted{\unsure{F}}F)&=\varepsilon_s(D)=\text{un des \inserted{deux} }n\text{-gones inscrits dans }\operatorname{Pol}(D)=\operatorname{Pol}(F),\\
\deleted{\varepsilon_{\unsure{s}}}\varepsilon_{\unsure{a}}(F)&=\varepsilon_a(F)=\text{\unsure{un|une|ens.} des deux }n\text{-gones circonscrits dans }\operatorname{Pol}(D)=\operatorname{Pol}(F)
\end{aligned} \]
